\chapter{Introduzione}
\label{cap:introduzione}

Questo progetto nasce nell'ambito del corso di \textit{Sistemi Operativi Dedicati} con l'obiettivo di realizzare un sistema embedded distribuito per il riconoscimento automatico delle targhe (\textit{Automatic License Plate Recognition}, ALPR), eseguito interamente su due microcontrollori comunicanti tra loro. Il progetto, denominato \textbf{ALPR distribuito}, è stato sviluppato su due schede STM32 Nucleo-H745ZI-Q e dimostra la fattibilità di una pipeline di intelligenza artificiale a due stadi: rilevamento e riconoscimento su hardware con risorse fortemente limitate, senza ricorrere ad acceleratori esterni o a un host più potente per l'inferenza. Il sistema riceve un'immagine da un PC, ne individua la regione contenente la targa attraverso una rete neurale di \textit{object detection}, ritaglia la regione di interesse (\textit{ROI}) e la trasmette a una seconda scheda che ne estrae il contenuto alfanumerico attraverso una rete di tipo \textit{transformer}.

\section{Motivazioni}
\label{sec:motivazioni-intro}
Negli ultimi anni si è affermato il paradigma dell'\textit{Edge AI}, in cui l'inferenza dei modelli neurali viene eseguita direttamente sul dispositivo di acquisizione invece che su un server remoto. Questo approccio è particolarmente interessante per un sistema ALPR, perché riduce la dipendenza dalla rete, contiene la latenza e limita la circolazione di dati sensibili: l'immagine può essere elaborata localmente e verso l'esterno viene comunicata solo la stringa della targa riconosciuta. Portare questa idea su microcontrollori STM32 permette inoltre di ottenere nodi a basso costo, con consumi ridotti e più semplici da installare rispetto a soluzioni basate su GPU o acceleratori embedded. In scenari con molti varchi o postazioni remote, queste caratteristiche incidono direttamente sulla sostenibilità economica ed energetica del sistema, rendendo possibile un dispiegamento più capillare anche in assenza di infrastrutture complesse.
Questi vantaggi sono, però, accompagnato da vincoli tecnici severi: la STM utilizzata, infatti, non integra unità dedicate per le reti neurali e affida l'intera inferenza al core Cortex-M7; di conseguenza, i modelli devono essere ridotti, quantizzati a 8 bit e organizzati con grande attenzione all'uso della memoria.

\begin{figure}[H]
    \centering
    \begin{tikzpicture}[
        font=\footnotesize,
        every node/.style={inner sep=2pt}
    ]
        % -------- Assi --------
        % Asse X: consumo (W) in scala log da 0.3 a 30
        % Asse Y: costo (EUR) in scala log da 10 a 1000
        % Mapping log -> coordinate cartesiane:
        %   X: 0.3 W -> 0  ;  3 W -> 4  ;  30 W -> 8
        %   Y: 10 EUR -> 0 ;  100 EUR -> 3 ;  1000 EUR -> 6
        % Quindi X in [0, 8], Y in [0, 6].
 
        % Rettangolo di sfondo
        \draw[->, thick] (0, 0) -- (8.5, 0) node[right] {Consumo [W]};
        \draw[->, thick] (0, 0) -- (0, 5.2) node[above] {Costo [EUR]};
 
        % Tick X (scala log)
        \foreach \x/\lab in {0/0.3, 2/1, 4/3, 6/10, 8/30}{
            \draw[thick] (\x, 0) -- (\x, -0.1) node[below] {\lab};
        }
 
        % Tick Y (scala log)
        \foreach \y/\lab in {0/10, 1.5/30, 3/100, 4.5/300}{
            \draw[thick] (0, \y) -- (-0.1, \y) node[left] {\lab};
        }
 
        % Griglia leggera
        \draw[gray!20, very thin] (0,0) grid[xstep=1, ystep=0.75] (8, 5);
 
        % -------- Punti (ognuno con label) --------
 
        % Nucleo STM32 H745  (0.5 W, 30 EUR)
        % log10(0.5/0.3) / log10(100) * 8  -> ~0.9 ; log10(30/10)/log10(100)*6 -> ~1.43
        \node[circle, fill=custom_red, minimum size=8pt, label={below right:\textbf{Nucleo H745}}] (a) at (0.9, 1.5) {};
 
        % Raspberry Pi 4  (5 W, 70 EUR)
        \node[circle, fill=custom_blue, minimum size=6pt, label={above:Raspberry Pi 4}] (b) at (4.9, 2.8) {};
 
        % Coral Dev Board  (3 W, 130 EUR)
        \node[circle, fill=custom_blue, minimum size=6pt, label={above:Coral Dev}] (c) at (4.0, 3.3) {};
 
        % Jetson Nano  (10 W, 150 EUR)
        \node[circle, fill=custom_blue, minimum size=6pt, label={right:Jetson Nano}] (d) at (6.0, 3.5) {};
 
        % Jetson Orin Nano  (15 W, 500 EUR)
        \node[circle, fill=custom_blue, minimum size=6pt, label={above:Jetson Orin}] (e) at (6.7, 4.6) {};
 
        % -------- Freccia "regione favorevole" --------
        \draw[->, thick, custom_red, dashed]
            (3.0, 4.5) to[bend right=20] node[midway, above, font=\scriptsize, custom_red, sloped] {edge AI estremo} (1.2, 1.8);
 
    \end{tikzpicture}
    \caption{Confronto costo-consumo tra la Nucleo H745 e altre piattaforme edge AI}
    \label{fig:edge-ai-comparison}
\end{figure}

\section{Obiettivi}
\label{sec:obiettivi}

Gli obiettivi principali del progetto sono i seguenti:

\begin{itemize}
    \item \textbf{Distribuire una pipeline ALPR su due microcontrollori}, affidando a ciascuna scheda un singolo stadio della pipeline (detection e OCR) e progettando un protocollo di comunicazione tra le due tramite UART.
    \item \textbf{Addestrare un detector custom}, basato sull'architettura YOLOv8 in una variante più piccola (denominata \textit{tiny}), su un dataset specifico per targhe automobilistiche, ottimizzandolo per immagini in scala di grigi a bassa risoluzione (192$\times$192~px) compatibili con i vincoli di memoria della scheda di esecuzione.
    \item \textbf{Integrare un modello OCR pre-addestrato} di tipo Compact Convolutional Transformer~\cite{hassani_2021_cct}, scelto per la sua capacità di operare con un numero di parametri ridotto, e adattarlo al formato di ingresso accettato dalla rete a partire dai dati prodotti dallo stadio di detection.
    \item \textbf{Sviluppare due firmware Zephyr} indipendenti che orchestrano comunicazione seriale, gestione dei buffer, esecuzione delle reti tramite il runtime \textbf{X-CUBE-AI} di \textit{STMicroelectronics} e sincronizzazione tra le due schede.
    \item \textbf{Misurare e documentare} le risorse impiegate (RAM, Flash) e le prestazioni del sistema reale, discutendo onestamente le limitazioni emerse e identificando i possibili sviluppi futuri.
\end{itemize}

\newpage
\section{Tecnologie utilizzate}
\label{sec:tecnologie}

Il sistema integra tecnologie eterogenee distribuite su tre livelli logici, ciascuno con un ruolo specifico:

\begin{itemize}
    \item \textbf{STM32 Nucleo-H745ZI-Q}: due schede di sviluppo dotate di un Cortex-M7 a 480~MHz, 1~MB di Flash on-chip e 512~KB di SRAM principale. La presenza di FPU hardware (FPv5-D16) è essenziale per gestire le operazioni in virgola mobile dei modelli neurali prima della quantizzazione.
    \item \textbf{Zephyr RTOS}: sistema operativo real-time open source utilizzato per il firmware di entrambe le schede. La sua architettura modulare basata su KConfig e devicetree consente di selezionare solo i moduli effettivamente necessari, riducendo l'occupazione di memoria e permettendo di rispettare i vincoli stringenti del progetto~\cite{zephyr_doc}.
    \item \textbf{X-CUBE-AI}: pacchetto di STMicroelectronics che converte modelli ONNX in librerie C ottimizzate per le serie STM32. Esegue la quantizzazione INT8 post-training e produce un runtime statico utilizzabile direttamente da Zephyr~\cite{stm32cubeai_doc}.
    \item \textbf{Ultralytics YOLOv8}~\cite{jocher_2023_yolov8}: framework PyTorch utilizzato per l'addestramento della rete di detection. La variante \textit{tiny} (scala \textit{t}) è stata personalizzata per accettare immagini monocromatiche e prevedere una sola classe.
    \item \textbf{fast-plate-ocr}~\cite{ankandrew_2024_fastplateocr}: repository open source che mette a disposizione il modello CCT-XS pre-addestrato per la lettura di targhe, esportato in formato ONNX e direttamente compatibile con il flusso di importazione di X-CUBE-AI.
    \item \textbf{Python 3 e OpenCV}: utilizzati lato host per pre-elaborare le immagini di test, gestire la comunicazione seriale tramite \texttt{pyserial} e visualizzare il risultato finale.
    \item \textbf{ST Edge AI Developer Cloud}: tool online messo a disposizione da \textit{STMicroelectronics} che consente di ottimizzare e convertire modelli di reti neurali per il deployment su microcontrollori STM32.
\end{itemize}

